\subsection{HMAC 消息认证码}

HMAC（Hash-based Message Authentication Code）是一种基于密码杂凑函数和共享密钥构造的消息认证码机制，由 Bellare、Canetti 和 Krawczyk 于 1996 年提出，并于 1997 年标准化为 RFC 2104\footnote{Bellare, M., Canetti, R., Krawczyk, H. "Keyed-Hashing for Message Authentication." RFC 2104, IETF, 1997.}。与单纯的哈希运算不同，HMAC 在计算过程中引入了密钥，因此不仅能够检测消息内容是否被篡改，还能够验证消息是否由持有合法密钥的一方生成。HMAC 通常用于身份认证、接口鉴权、请求防伪造和数据完整性校验等场景。

设 $H$ 表示密码杂凑函数，$K$ 表示共享密钥，$m$ 表示待认证消息，则 HMAC 的一般形式可以表示为：
\begin{equation}
\operatorname{HMAC}(K,m)=H\bigl((K' \oplus \mathit{opad})\parallel H((K' \oplus \mathit{ipad})\parallel m)\bigr)
\end{equation}
其中，$K'$ 表示经过填充或压缩处理后的密钥（若密钥长度小于杂凑函数的分组长度，则在右侧填充零；若大于分组长度，则先对密钥计算一次杂凑值），$\mathit{opad}$（outer pad）和 $\mathit{ipad}$（inner pad）分别表示外部填充值和内部填充值，$\parallel$ 表示字符串连接操作，$\oplus$ 表示按位异或。对于输出 256 bit 摘要的杂凑函数（如 SM3），$\mathit{ipad}$ 为字节 0x36 重复 32 次构成的 512 bit 序列，$\mathit{opad}$ 为字节 0x5C 重复 32 次构成的 512 bit 序列。

HMAC 的安全性取决于底层杂凑函数的性质和密钥的保密性。在形式化安全分析中，HMAC 的安全性通常以\textbf{选择消息攻击下的存在性不可伪造性}（Existential Unforgeability under Chosen-Message Attack, EUF-CMA）来刻画：对于概率多项式时间的攻击者 $\mathcal{A}$，若其不知道密钥 $K$，但在可以查询任意消息 $m_i$ 对应的 $\operatorname{HMAC}(K, m_i)$ 的条件下，仍然难以构造出一个未查询过的消息 $m^*$ 的合法认证码 $\operatorname{HMAC}(K, m^*)$，则称该 HMAC 方案满足 EUF-CMA 安全性。当底层杂凑函数满足抗碰撞性和伪随机函数性质时，HMAC 的安全性可以归约到杂凑函数的安全性质上。

% !TEX root = ../../main.tex
% HMAC 结构图 —— 双层哈希流向
\begin{figure}[H]
\centering
\begin{tikzpicture}[
  node distance = 0.5cm and 0.7cm,
  box/.style    = {rectangle, draw, rounded corners=3pt, minimum width=1.8cm, minimum height=0.6cm,
                   font=\small, fill=white},
  op/.style     = {circle, draw, minimum size=0.5cm, font=\small\bfseries, fill=gray!15},
  arr/.style    = {-Stealth, thick},
  lbl/.style    = {font=\small\itshape},
  grn/.style    = {fill=green!12, draw=green!50!black},
  blu/.style    = {fill=blue!10, draw=blue!50!black},
  orn/.style    = {fill=orange!18, draw=orange!70!black},
]

% 原始密钥 K
\node[box, blu] (K) {密钥$K$};

% 密钥预处理
\node[box, blu, below=0.65cm of K] (Kp) {$K'$（填充/截断至分组长度）};

% ipad / opad（左右对称）
\node[box, orn, below left=0.8cm and 0.4cm of Kp]  (ipad) {$\mathit{ipad}$};
\node[box, orn, below right=0.8cm and 0.4cm of Kp] (opad) {$\mathit{opad}$};

% XOR nodes（在 ipad/opad 下方）
\node[op, below=0.55cm of ipad] (xori) {$\oplus$};
\node[op, below=0.55cm of opad] (xoro) {$\oplus$};

% 消息 m（左侧，在 xori 旁边）
\node[box, blu, left=0.7cm of xori] (msg) {消息$m$};

% 内层哈希输入拼接
\node[box, grn, below=0.5cm of xori, minimum width=3.0cm] (inner_in) {$(K' \oplus \mathit{ipad}) \;\|\; m$};

% 内层哈希
\node[box, grn, below=0.45cm of inner_in] (H1) {$H(\cdot)$（内层哈希）};

% 外层拼接
\node[box, orn, below=0.6cm of H1, minimum width=3.0cm] (outer_in) {$(K' \oplus \mathit{opad}) \;\|\; H(\text{inner})$};

% 外层哈希
\node[box, orn, below=0.45cm of outer_in] (H2) {$H(\cdot)$（外层哈希）};

% 输出
\node[box, below=0.45cm of H2, minimum width=2.4cm, fill=red!10, draw=red!50!black] (out) {$\operatorname{HMAC}(K, m)$};

% 箭头：K → Kp
\draw[arr] (K) -- (Kp);

% 箭头：Kp → ipad, Kp → opad（分支）
\draw[arr] (Kp.south west) -- (ipad.north);
\draw[arr] (Kp.south east) -- (opad.north);

% 箭头：Kp → xori（从 Kp 左侧绕行，不穿越 ipad）
\draw[arr] (Kp.west) -- ++(-0.35,0) |- (xori.north);

% 箭头：Kp → xoro（从 Kp 右侧绕行，不穿越 opad）
\draw[arr] (Kp.east) -- ++(0.35,0) |- (xoro.north);

% 箭头：ipad → xori, opad → xoro
\draw[arr] (ipad) -- (xori);
\draw[arr] (opad) -- (xoro);

% 箭头：msg → xori
\draw[arr] (msg) -- (xori);

% 箭头：内层流程
\draw[arr] (xori) -- (inner_in);
\draw[arr] (inner_in) -- (H1);
\draw[arr] (H1) -- (outer_in);

% 箭头：xoro → outer_in（从右侧绕行）
\draw[arr] (xoro.south) -- ++(0,-0.25) -| (outer_in.north east);

% 箭头：外层流程
\draw[arr] (outer_in) -- (H2);
\draw[arr] (H2) -- (out);

% 注释
\node[lbl, right=0.15cm of xori] {内层混合};
\node[lbl, right=0.15cm of H2]   {外层混合};

\end{tikzpicture}
\caption{HMAC 双层嵌套哈希结构}
\label{fig:hmac_structure}
\end{figure}


在 FoodShield 系统中，HMAC 主要用于订单 Token 的生成与验证。用户在创建订单后，平台服务器根据订单号、匿名身份标识、时间戳和随机数等字段生成订单认证 Token。用户进入订单通信通道时需要提交订单号和 Token，服务器重新计算并比对 Token，以判断用户是否为该订单的合法参与方。该机制使用户能够在不直接暴露真实身份的情况下证明自己拥有订单通信资格。

本系统中订单 Token 的抽象形式如下：

\begin{equation}
m_{\mathrm{order}} = \mathit{orderID} \parallel \mathit{PID} \parallel \mathit{timestamp} \parallel \mathit{nonce}
\end{equation}

\begin{equation}
\mathrm{token} = \operatorname{HMAC}(K_{\mathrm{order}},\, m_{\mathrm{order}})
\end{equation}

其中，$K_{\mathrm{order}}$为订单认证密钥，$\mathit{orderID}$为订单编号，$\mathit{PID}$为用户匿名身份标识，$\mathit{timestamp}$为时间戳，$\mathit{nonce}$为 128 bit 随机数。通过引入时间戳和随机数，可以降低 Token 被重放利用的风险——即使攻击者截获了某次通信的 Token，由于时间戳的时效性约束和随机数的不可预测性，该 Token 也无法在其他订单或其他时间窗口中被成功复用。通过绑定订单号与匿名身份标识，可以防止攻击者将一个订单的认证凭证迁移到其他订单中使用。
